Action-conditional predictive objectives are a standard way to train history representations in partially observable environments, and a resettable simulator makes it tempting to supervise them with the outcomes of several alternative actions taken from the same saved state. Such branching consumes environment transitions and relies on a privilege, state restore, that ordinary interaction does not provide. We study how to account for this privilege when comparing representation learners. We implement a small aliased T-maze in which the current observation at a junction is identical for both values of an earlier clue, a transition ledger that charges every main-trajectory, branch and evaluation step against one hard cap and counts restores separately, and a fixed readout protocol that trains a fresh action-conditional reward head on a frozen encoder and scores a model-predictive planner by its junction decision (forced-commit accuracy). An evaluator diagnostic with fixed encoders shows that this readout can separate informative from uninformative histories: a hand-designed history-only control reaches forced-commit accuracy \DiagForcedInfoMin{} in all three diagnostic seeds, while an untrained recurrent encoder and a memoryless control both stay at \DiagForcedRandMax{}, leaving a headroom of \DiagHeadroomMin{}. In the same runs, the full-action planning success of the two clue-blind encoders differs only through stalling behaviour, so we do not use it to compare representations. We then ran a 2$\times$2 study that crosses action conditioning (an action-conditional predictor versus an action-free ablation) with the collection method (branching versus no branching). Each collection method had \StudyCfgPretrainBudget{} pretraining transitions per seed, and every encoder was read out through a separate restore-free pool. The readout again recognised the informative control, at forced-commit accuracy \StudyForcedInfoMin{} in all seeds. None of the four learned encoders exceeded the untrained encoder: their accuracies ranged from \StudyForcedLearnedMin{} to \StudyForcedLearnedMax{}, so the pre-registered learning-benefit contrast failed (\StudyVerdict). Because every learned arm stayed at chance, the collection and conditioning contrasts sit at a floor and do not separate the two effects; we do not read them as evidence of no effect.We make no claim that low prediction error guarantees planning for arbitrary policies, and all results come from one synthetic task family with three seeds.
